\begin{tikzpicture}[
  kotak/.style={draw, rounded corners, align=center, text width=3.3cm,
    minimum height=1.3cm, font=\scriptsize, fill=blue!8},
  masalah/.style={draw, rounded corners, align=center, text width=3.3cm,
    minimum height=0.75cm, font=\scriptsize, fill=red!22},
  panah/.style={-{Stealth[length=2mm]}, thick},
  keterangan/.style={font=\tiny, align=center, inner sep=1pt}
]
  \node[masalah] (masalah1) at (0, 4.0) {Bila hilang:\\keputusan tanpa pemutus};
  \node[masalah] (masalah2) at (4.0, 4.0) {Bila hilang:\\keputusan tidak ditinjau};
  \node[masalah] (masalah3) at (8.0, 4.0) {Bila hilang:\\piagam hanya tersimpan};
  \node[kotak] (struktur) at (0, 2.1)
    {\textbf{Struktur}\\{\tiny\mdseries forum, peran, mandat,\\dan batas delegasi}};
  \node[kotak] (proses) at (4.0, 2.1)
    {\textbf{Proses}\\{\tiny\mdseries agenda EDM, catatan keputusan,\\dan peninjauan
    manfaat}};
  \node[kotak] (relasional) at (8.0, 2.1)
    {\textbf{Mekanisme relasional}\\{\tiny\mdseries dialog lintas fungsi, bahasa
    bersama,\\dan belajar dari insiden}};
  \node[draw, rounded corners, align=center, fill=green!10, font=\scriptsize,
    text width=11.1cm, minimum height=0.9cm] (hasil) at (4.0, 0)
    {\textbf{Keputusan yang dijalankan dan dapat diperiksa}\\{\tiny\mdseries satu
    pemutus, alasan tercatat, tindak lanjut selesai, dan hasil layanan ditinjau}};
  \draw[panah, dashed] (struktur) -- (masalah1);
  \draw[panah, dashed] (proses) -- (masalah2);
  \draw[panah, dashed] (relasional) -- (masalah3);
  \draw[panah] (struktur) -- (struktur |- hasil.north);
  \draw[panah] (proses) -- (proses |- hasil.north);
  \draw[panah] (relasional) -- (relasional |- hasil.north);
\end{tikzpicture}
